\begin{tikzpicture}[
  font=\sffamily\scriptsize,
  block/.style={
    draw=ctline, rounded corners=2mm, thick, fill=white,
    minimum height=13mm, align=center
  },
  physical/.style={block, draw=ctblue, fill=blue!5},
  logical/.style={block, draw=ctgreen, fill=ctmint},
  pattern/.style={block, draw=ctamber, fill=yellow!7},
  flow/.style={-{Latex[length=2mm]}, thick, draw=ctgray},
  note/.style={font=\sffamily\tiny, text=ctgray, align=center}
]
  \node[physical, text width=38mm] (client) at (-4.6,3.1)
    {\textbf{Client web}\\application React};
  \node[physical, text width=38mm] (server) at (0,3.1)
    {\textbf{Serveur applicatif}\\API Laravel};
  \node[physical, text width=38mm] (data) at (4.6,3.1)
    {\textbf{Serveur de données}\\PostgreSQL};

  \draw[flow] (client) -- node[note, above]{HTTPS / JSON} (server);
  \draw[flow] (server) -- node[note, above]{SQL transactionnel} (data);

  \node[logical, text width=35mm] (presentation) at (-4.6,0.6)
    {\textbf{Présentation}\\pages, composants, état UI};
  \node[logical, text width=35mm] (application) at (0,0.6)
    {\textbf{Application et métier}\\cas d'utilisation, règles};
  \node[logical, text width=35mm] (persistence) at (4.6,0.6)
    {\textbf{Persistance}\\modèles et accès aux données};

  \draw[flow] (presentation) -- (application);
  \draw[flow] (application) -- (persistence);

  \node[pattern, text width=35mm] (components) at (-4.6,-2.0)
    {\textbf{Composants React}\\composition fonctionnelle};
  \node[pattern, text width=35mm] (mvc) at (0,-2.0)
    {\textbf{MVC côté Laravel}\\contrôleur, modèle, réponse API};
  \node[pattern, text width=35mm] (boundaries) at (4.6,-2.0)
    {\textbf{Adaptateurs et événements}\\intégrations remplaçables};

  \draw[flow] (presentation) -- (components);
  \draw[flow] (application) -- (mvc);
  \draw[flow] (persistence) -- (boundaries);

  \node[note, text width=125mm] at (0,-3.35)
    {Architecture client--serveur et trois tiers, réalisée sous la forme d'un
    monolithe modulaire complété par une passerelle OCPP spécialisée.};
\end{tikzpicture}
